% Options for packages loaded elsewhere
\PassOptionsToPackage{unicode}{hyperref}
\PassOptionsToPackage{hyphens}{url}
\PassOptionsToPackage{dvipsnames,svgnames,x11names}{xcolor}
%
\documentclass[
  10pt,
]{article}
\usepackage{amsmath,amssymb}
\usepackage{iftex}
\ifPDFTeX
  \usepackage[T1]{fontenc}
  \usepackage[utf8]{inputenc}
  \usepackage{textcomp} % provide euro and other symbols
\else % if luatex or xetex
  \usepackage{unicode-math} % this also loads fontspec
  \defaultfontfeatures{Scale=MatchLowercase}
  \defaultfontfeatures[\rmfamily]{Ligatures=TeX,Scale=1}
\fi
\usepackage{lmodern}
\ifPDFTeX\else
  % xetex/luatex font selection
  \setmainfont[]{DejaVu Serif}
  \setsansfont[]{DejaVu Sans}
  \setmonofont[]{DejaVu Sans Mono}
  \setmathfont[]{Latin Modern Math}
\fi
% Use upquote if available, for straight quotes in verbatim environments
\IfFileExists{upquote.sty}{\usepackage{upquote}}{}
\IfFileExists{microtype.sty}{% use microtype if available
  \usepackage[]{microtype}
  \UseMicrotypeSet[protrusion]{basicmath} % disable protrusion for tt fonts
}{}
\makeatletter
\@ifundefined{KOMAClassName}{% if non-KOMA class
  \IfFileExists{parskip.sty}{%
    \usepackage{parskip}
  }{% else
    \setlength{\parindent}{0pt}
    \setlength{\parskip}{6pt plus 2pt minus 1pt}}
}{% if KOMA class
  \KOMAoptions{parskip=half}}
\makeatother
\usepackage{xcolor}
\usepackage[margin=2.0cm]{geometry}
\usepackage{color}
\usepackage{fancyvrb}
\newcommand{\VerbBar}{|}
\newcommand{\VERB}{\Verb[commandchars=\\\{\}]}
\DefineVerbatimEnvironment{Highlighting}{Verbatim}{commandchars=\\\{\}}
% Add ',fontsize=\small' for more characters per line
\newenvironment{Shaded}{}{}
\newcommand{\AlertTok}[1]{\textcolor[rgb]{1.00,0.00,0.00}{\textbf{#1}}}
\newcommand{\AnnotationTok}[1]{\textcolor[rgb]{0.38,0.63,0.69}{\textbf{\textit{#1}}}}
\newcommand{\AttributeTok}[1]{\textcolor[rgb]{0.49,0.56,0.16}{#1}}
\newcommand{\BaseNTok}[1]{\textcolor[rgb]{0.25,0.63,0.44}{#1}}
\newcommand{\BuiltInTok}[1]{\textcolor[rgb]{0.00,0.50,0.00}{#1}}
\newcommand{\CharTok}[1]{\textcolor[rgb]{0.25,0.44,0.63}{#1}}
\newcommand{\CommentTok}[1]{\textcolor[rgb]{0.38,0.63,0.69}{\textit{#1}}}
\newcommand{\CommentVarTok}[1]{\textcolor[rgb]{0.38,0.63,0.69}{\textbf{\textit{#1}}}}
\newcommand{\ConstantTok}[1]{\textcolor[rgb]{0.53,0.00,0.00}{#1}}
\newcommand{\ControlFlowTok}[1]{\textcolor[rgb]{0.00,0.44,0.13}{\textbf{#1}}}
\newcommand{\DataTypeTok}[1]{\textcolor[rgb]{0.56,0.13,0.00}{#1}}
\newcommand{\DecValTok}[1]{\textcolor[rgb]{0.25,0.63,0.44}{#1}}
\newcommand{\DocumentationTok}[1]{\textcolor[rgb]{0.73,0.13,0.13}{\textit{#1}}}
\newcommand{\ErrorTok}[1]{\textcolor[rgb]{1.00,0.00,0.00}{\textbf{#1}}}
\newcommand{\ExtensionTok}[1]{#1}
\newcommand{\FloatTok}[1]{\textcolor[rgb]{0.25,0.63,0.44}{#1}}
\newcommand{\FunctionTok}[1]{\textcolor[rgb]{0.02,0.16,0.49}{#1}}
\newcommand{\ImportTok}[1]{\textcolor[rgb]{0.00,0.50,0.00}{\textbf{#1}}}
\newcommand{\InformationTok}[1]{\textcolor[rgb]{0.38,0.63,0.69}{\textbf{\textit{#1}}}}
\newcommand{\KeywordTok}[1]{\textcolor[rgb]{0.00,0.44,0.13}{\textbf{#1}}}
\newcommand{\NormalTok}[1]{#1}
\newcommand{\OperatorTok}[1]{\textcolor[rgb]{0.40,0.40,0.40}{#1}}
\newcommand{\OtherTok}[1]{\textcolor[rgb]{0.00,0.44,0.13}{#1}}
\newcommand{\PreprocessorTok}[1]{\textcolor[rgb]{0.74,0.48,0.00}{#1}}
\newcommand{\RegionMarkerTok}[1]{#1}
\newcommand{\SpecialCharTok}[1]{\textcolor[rgb]{0.25,0.44,0.63}{#1}}
\newcommand{\SpecialStringTok}[1]{\textcolor[rgb]{0.73,0.40,0.53}{#1}}
\newcommand{\StringTok}[1]{\textcolor[rgb]{0.25,0.44,0.63}{#1}}
\newcommand{\VariableTok}[1]{\textcolor[rgb]{0.10,0.09,0.49}{#1}}
\newcommand{\VerbatimStringTok}[1]{\textcolor[rgb]{0.25,0.44,0.63}{#1}}
\newcommand{\WarningTok}[1]{\textcolor[rgb]{0.38,0.63,0.69}{\textbf{\textit{#1}}}}
\usepackage{longtable,booktabs,array}
\usepackage{calc} % for calculating minipage widths
% Correct order of tables after \paragraph or \subparagraph
\usepackage{etoolbox}
\makeatletter
\patchcmd\longtable{\par}{\if@noskipsec\mbox{}\fi\par}{}{}
\makeatother
% Allow footnotes in longtable head/foot
\IfFileExists{footnotehyper.sty}{\usepackage{footnotehyper}}{\usepackage{footnote}}
\makesavenoteenv{longtable}
\setlength{\emergencystretch}{3em} % prevent overfull lines
\providecommand{\tightlist}{%
  \setlength{\itemsep}{0pt}\setlength{\parskip}{0pt}}
\setcounter{secnumdepth}{-\maxdimen} % remove section numbering
% doc_header.tex -- LaTeX preamble additions used by build_docs.sh for every
% Markdown -> LaTeX -> PDF conversion (supplements, README, submission files).
\usepackage{etoolbox}
% longtable may not be loaded by pandoc when a document has no table
\usepackage{longtable}
% bold vectors: \bm -> \symbf under unicode-math (XeLaTeX), bm package otherwise
\makeatletter
\@ifpackageloaded{unicode-math}{\providecommand{\bm}[1]{\symbf{#1}}}{\usepackage{bm}}
\makeatother
% smaller type in long tables so that wide audit tables fit the text width;
% inside a table a line break is allowed after an underscore of plain text
% (pandoc writes it as \_), so that long identifiers wrap inside their column
% rows are set a little apart, so that symbols with indices do not touch
\AtBeginEnvironment{longtable}{\footnotesize\renewcommand{\arraystretch}{1.45}%
  \let\sptunderscore\_\def\_{\sptunderscore\allowbreak}}
\setlength{\LTcapwidth}{\textwidth}
% allow line breaks after slashes and hyphens inside table cells
\usepackage{xurl}
% code blocks: break over-long lines instead of running into the margin.
% Plain ``` blocks are typeset by pandoc as {verbatim}, highlighted ones
% (```bash) as {Highlighting}; both are re-declared as breakable fancyvrb
% environments.  The continuation line is marked with a small hooked arrow.
\usepackage{fvextra}
\RecustomVerbatimEnvironment{verbatim}{Verbatim}{breaklines,breakanywhere,fontsize=\small}
\makeatletter
\@ifundefined{Highlighting}{}{%
  \RecustomVerbatimEnvironment{Highlighting}{Verbatim}{breaklines,breakanywhere,fontsize=\small,commandchars=\\\{\}}}
\makeatother
% inline code (pandoc writes it as \texttt{...}): never hyphenate -- an inserted
% hyphen could be read as part of a file name -- but allow a line break after
% an underscore, slash, comma, colon or equals sign, so that long paths and
% column lists wrap inside the text width and inside table columns.
% Inside a formula (e.g. the name of a notebook symbol in an equation) the same
% command must first leave math mode, which \text of amsmath does.
\ExplSyntaxOn
\tl_new:N \l__spt_code_tl
\cs_new_protected:Npn \__spt_texttt:n #1
  {
    \tl_set:Nn \l__spt_code_tl {#1}
    \regex_replace_all:nnN { ([\/\,\:\=]) } { \1 \c{allowbreak} } \l__spt_code_tl
    \regex_replace_all:nnN { \c{_} } { \c{_} \c{allowbreak} } \l__spt_code_tl
    \group_begin: \ttfamily \hyphenchar \font = -1 \scan_stop: \l__spt_code_tl \group_end:
  }
\RenewDocumentCommand \texttt { m }
  {
    \mode_if_math:TF { \text { \__spt_texttt:n {#1} } } { \__spt_texttt:n {#1} }
  }
\ExplSyntaxOff
% A figure is written in the Markdown source as an image paragraph followed by
% a caption paragraph ("**Figure S1.** ..."), so that the caption is also
% visible on GitHub.  An image that starts a paragraph is centered.
\usepackage{graphicx}% pandoc loads it only for documents that contain an image
\NewCommandCopy\sptincludegraphics\includegraphics
\RenewDocumentCommand\includegraphics{O{} m}{%
  \ifvmode\begin{center}\sptincludegraphics[#1]{#2}\end{center}%
  \else\sptincludegraphics[#1]{#2}\fi}
% keep figures where they are referenced
\usepackage{float}
\floatplacement{figure}{H}
% running header with the document title
\usepackage{fancyhdr}
\pagestyle{fancy}
\fancyhf{}
\fancyhead[L]{\small PROJECT spacetimes, Revision\_Deflation\_CJT\_and\_Kohn-Sham}
\fancyhead[R]{\small\thepage}
\renewcommand{\headrulewidth}{0.4pt}
% 16x16 matrices: tighter columns so that a full matrix fits the text width
\setlength{\arraycolsep}{3.2pt}
% The documents of this project contain long tagged symbols, file paths and
% identifiers that cannot be hyphenated.  A paragraph is therefore allowed to
% have looser lines before TeX gives up and lets a line stick out of the margin.
\AtBeginDocument{\tolerance=2000 \emergencystretch=5em \hbadness=2000}
\ifLuaTeX
  \usepackage{selnolig}  % disable illegal ligatures
\fi
\IfFileExists{bookmark.sty}{\usepackage{bookmark}}{\usepackage{hyperref}}
\IfFileExists{xurl.sty}{\usepackage{xurl}}{} % add URL line breaks if available
\urlstyle{same}
\hypersetup{
  pdftitle={README --- Revision\_Deflation\_CJT\_and\_Kohn-Sham},
  pdfauthor={Patrick Lee Nash, Ph.D.},
  colorlinks=true,
  linkcolor={blue},
  filecolor={Maroon},
  citecolor={blue},
  urlcolor={blue},
  pdfcreator={LaTeX via pandoc}}

\title{README --- Revision\_Deflation\_CJT\_and\_Kohn-Sham}
\usepackage{etoolbox}
\makeatletter
\providecommand{\subtitle}[1]{% add subtitle to \maketitle
  \apptocmd{\@title}{\par {\large #1 \par}}{}{}
}
\makeatother
\subtitle{PROJECT spacetimes: the 16-component Dirac field with mass
term and self-interaction on the 4+4 and 7+1 spacetimes whose
coordinates x5, x6, x7 deflate exponentially; Kohn--Sham and CJT
analyses}
\author{Patrick Lee Nash, Ph.D.}
\date{November 13, 2026}

\begin{document}
\maketitle

\hypertarget{what-this-folder-is}{%
\section{1. What this folder is}\label{what-this-folder-is}}

The folder holds one sub-project of PROJECT spacetimes. Its hypotheses
are the author's:

\begin{itemize}
\tightlist
\item
  The three extra coordinates \(x5\), \(x6\), \(x7\) deflate
  exponentially in the ordinary time \(x4\). The metric is diagonal, and
  its components along them are
  \(g^{(4+4)}_{tt}=-\mathrm{const}_t^{\,2}\,e^{-2\,\mathrm{rate}_t\,x4}\)
  on the 4+4 spacetime (where they are time-like) and
  \(g^{(7+1)}_{tt}=+\mathrm{const}_t^{\,2}\,e^{-2\,\mathrm{rate}_t\,x4}\)
  on the 7+1 spacetime (where they are space-like), for \(t=5,6,7\),
  with the positive rates \(\mathrm{rate}_5=\mathrm{g55}\),
  \(\mathrm{rate}_6=\mathrm{g66}\), \(\mathrm{rate}_7=\mathrm{g77}\) and
  positive constants \(\mathrm{const}_t\).
\item
  The Lagrangian of the sixteen-component Dirac field contains a kinetic
  term, a \textbf{mass term} and a \textbf{self-interaction} that is
  quartic in the field, with nine couplings \(g_0,\dots,g_8\).
\item
  The consequences are worked out with two methods, the Kohn--Sham
  ansatz and the effective potential for composite operators of
  Cornwall, Jackiw and Tomboulis (CJT), each as a check of the other and
  of the earlier results.
\end{itemize}

Everything in the folder was computed for it: every number in a document
comes from a data file that a command of the program wrote, and every
table is inserted by a script.

\hypertarget{where-to-start-reading}{%
\section{2. Where to start reading}\label{where-to-start-reading}}

\begin{longtable}[]{@{}
  >{\raggedright\arraybackslash}p{(\columnwidth - 2\tabcolsep) * \real{0.5000}}
  >{\raggedright\arraybackslash}p{(\columnwidth - 2\tabcolsep) * \real{0.5000}}@{}}
\toprule\noalign{}
\begin{minipage}[b]{\linewidth}\raggedright
if you want
\end{minipage} & \begin{minipage}[b]{\linewidth}\raggedright
read
\end{minipage} \\
\midrule\noalign{}
\endhead
\bottomrule\noalign{}
\endlastfoot
the physics of one spacetime, complete & the manuscript for that
spacetime, in \texttt{manuscript/}; it defines every symbol it uses and
contains its own appendices \\
the definitions: metric, Dirac matrices printed in full, Lagrangians,
field equations, derived step by step for both spacetimes side by side &
\texttt{dirac\ matrices\ deflating} \\
the meaning of a symbol &
\texttt{provenance/P1\_master\_symbols\_spacetimes} (every symbol of the
project, with one meaning) \\
how a number was computed, or how to compute it again &
\texttt{provenance/P2\_calculations\_code\_and\_checks} \\
what was asked, what was done, what was not, and which decisions are
open & \texttt{provenance/P0\_task\_rules\_deviations\_decisions} \\
\end{longtable}

Each document exists as \texttt{.md}, \texttt{.tex} and \texttt{.pdf},
and each is complete in itself.

\hypertarget{contents}{%
\section{3. Contents}\label{contents}}

\begin{longtable}[]{@{}
  >{\raggedright\arraybackslash}p{(\columnwidth - 2\tabcolsep) * \real{0.5000}}
  >{\raggedright\arraybackslash}p{(\columnwidth - 2\tabcolsep) * \real{0.5000}}@{}}
\toprule\noalign{}
\begin{minipage}[b]{\linewidth}\raggedright
path
\end{minipage} & \begin{minipage}[b]{\linewidth}\raggedright
content
\end{minipage} \\
\midrule\noalign{}
\endhead
\bottomrule\noalign{}
\endlastfoot
\texttt{manuscript/16\ Component\ Dirac\ Fermion\ Condensate\ Singlet\ Superfluid\ on\ 4+4\ Deflating\ Spacetime}
& manuscript for the 4+4 spacetime (\texttt{.tex}, \texttt{.pdf},
\texttt{.md}) \\
\texttt{manuscript/16\ Component\ Dirac\ Fermion\ Condensate\ Singlet\ Superfluid\ on\ 7+1\ Deflating\ Spacetime}
& manuscript for the 7+1 spacetime \\
\texttt{manuscript/deflation\_refs.bib} & the references of the
manuscripts \\
\texttt{dirac\ matrices\ deflating} & geometry, Dirac matrices,
Lagrangians, field equations, solutions; reports of the checks \\
\texttt{provenance/P0\_task\_rules\_deviations\_decisions} & the task,
the rules, deviations, decisions, error log \\
\texttt{provenance/P1\_master\_symbols\_spacetimes} & the MASTER
registry of symbols and the proof that no symbol has two meanings \\
\texttt{provenance/P2\_calculations\_code\_and\_checks} & commands,
instructions, data files, figures, complete output of all checks, source
files \\
\texttt{provenance/symbols\_master.csv} & the registry as a table (229
symbols) \\
\texttt{provenance/generated/} & tables and numbers written by the
scripts from the data files \\
\texttt{code/deflation-cjt-ks-rs/} & the Rust program
\texttt{deflation\_cjt\_ks} (33 source files, 15007 lines, 68 unit
tests) \\
\texttt{code/scripts/} & independent verification, figures, and the
scripts that build and check the documents \\
\texttt{data/} & 111 data files and the reports of the checks \\
\texttt{figures/} & 30 figures (\texttt{.pdf} and \texttt{.png}) \\
\texttt{docs\_src/} & the sources of all documents \\
\end{longtable}

\hypertarget{how-to-run-everything}{%
\section{4. How to run everything}\label{how-to-run-everything}}

\textbf{What is needed.} A Rust toolchain (\texttt{cargo}); Python 3
with numpy, scipy, sympy, mpmath and matplotlib; for the documents,
pandoc, XeLaTeX with latexmk, and pdflatex with REVTeX 4.2 and BibTeX.
The program uses the Rust crate \texttt{code/spacetimes-rs} at the top
of the repository as a library; nothing is downloaded.

\textbf{Where.} All commands are run from this folder. On the author's
computer:

\begin{Shaded}
\begin{Highlighting}[]
\BuiltInTok{cd} \StringTok{"/Users/nsh/Developer/github/spacetimes/Revision\_Deflation\_CJT\_and\_Kohn{-}Sham"}
\end{Highlighting}
\end{Shaded}

\textbf{Step 1. Compile the program and run its unit tests.}

\begin{Shaded}
\begin{Highlighting}[]
\ExtensionTok{cargo}\NormalTok{ test  }\AttributeTok{{-}{-}release} \AttributeTok{{-}{-}manifest{-}path}\NormalTok{ code/deflation{-}cjt{-}ks{-}rs/Cargo.toml}
\ExtensionTok{cargo}\NormalTok{ build }\AttributeTok{{-}{-}release} \AttributeTok{{-}{-}manifest{-}path}\NormalTok{ code/deflation{-}cjt{-}ks{-}rs/Cargo.toml}
\end{Highlighting}
\end{Shaded}

The first command must end with \texttt{test\ result:\ ok}. The second
builds the program
\texttt{code/deflation-cjt-ks-rs/target/release/deflation\_cjt\_ks}.

\textbf{Step 2. Run the commands of the program.} Each writes its data
files and its report of checks into \texttt{data} and ends with a line
such as \texttt{kohnsham:\ 13\ of\ 13\ checks\ passed}.

\begin{Shaded}
\begin{Highlighting}[]
\VariableTok{B}\OperatorTok{=}\NormalTok{code/deflation{-}cjt{-}ks{-}rs/target/release/deflation\_cjt\_ks}
\VariableTok{$B}\NormalTok{ help}
\ControlFlowTok{for}\NormalTok{ c }\KeywordTok{in}\NormalTok{ selfcoupling kohnsham pairing gravity matrices modes symmetry channels }\DataTypeTok{\textbackslash{}}
\NormalTok{         history quench paired cosmology matter background spectrum}\KeywordTok{;} \ControlFlowTok{do} \VariableTok{$B} \VariableTok{$c}\KeywordTok{;} \ControlFlowTok{done}
\VariableTok{$B}\NormalTok{ forecast}
\VariableTok{$B}\NormalTok{ superfluid}
\VariableTok{$B}\NormalTok{ families}
\end{Highlighting}
\end{Shaded}

The first line abbreviates the path of the program, and \texttt{help}
lists the commands with a description of each. The loop runs fifteen
commands, each of which takes from a few seconds to about two minutes.
The last three take longer (about a quarter of an hour, three quarters
of an hour, and an hour and a quarter on a machine with two processors).

\textbf{Step 3. The independent verifications.}

\begin{Shaded}
\begin{Highlighting}[]
\ExtensionTok{python3}\NormalTok{ code/scripts/verify\_deflating\_independent.py}
\ExtensionTok{python3}\NormalTok{ code/scripts/verify\_cjt\_ks\_independent.py}
\ExtensionTok{python3}\NormalTok{ code/scripts/mutation\_test\_independent.py}
\ExtensionTok{python3}\NormalTok{ code/scripts/vendor\_eb.py }\AttributeTok{{-}{-}check}
\end{Highlighting}
\end{Shaded}

The first two repeat the central results with computer algebra and
general-purpose numerical tools (43 of 43 and 37 of 37 checks pass). The
third introduces deliberate errors into copies of the two scripts, one
at a time, and records whether the checks notice (36 of 36 detected;
about six minutes). The fourth verifies that four modules of the program
are exact copies of modules of the library crate with a listed set of
changes.

\textbf{Step 4. Figures, tables, documents, and the checks of the
documents.}

\begin{Shaded}
\begin{Highlighting}[]
\ExtensionTok{python3}\NormalTok{ code/scripts/make\_figures.py}
\ExtensionTok{python3}\NormalTok{ code/scripts/write\_symbol\_registry.py}
\ExtensionTok{python3}\NormalTok{ code/scripts/write\_doc\_tables.py}
\ExtensionTok{python3}\NormalTok{ code/scripts/assemble\_docs.py}
\ExtensionTok{python3}\NormalTok{ code/scripts/build\_manuscripts.py}
\FunctionTok{bash}\NormalTok{    code/scripts/build\_docs.sh}
\FunctionTok{bash}\NormalTok{    code/scripts/check\_all.sh}
\end{Highlighting}
\end{Shaded}

In this order: draw every figure from the data files; write the registry
of symbols; write every table and every computed number of the documents
from the data files; assemble the Markdown documents; assemble and
typeset the two manuscripts; typeset the other documents; run the checks
of the documents and read the verdict of every report. The last command
prints one line for each check and must end with \texttt{ALL\ PASS}.

\hypertarget{the-checks-in-numbers}{%
\section{5. The checks, in numbers}\label{the-checks-in-numbers}}

\begin{longtable}[]{@{}
  >{\raggedright\arraybackslash}p{(\columnwidth - 2\tabcolsep) * \real{0.5000}}
  >{\raggedright\arraybackslash}p{(\columnwidth - 2\tabcolsep) * \real{0.5000}}@{}}
\toprule\noalign{}
\begin{minipage}[b]{\linewidth}\raggedright
what
\end{minipage} & \begin{minipage}[b]{\linewidth}\raggedright
result
\end{minipage} \\
\midrule\noalign{}
\endhead
\bottomrule\noalign{}
\endlastfoot
checks of the program (exact where possible; reports
\texttt{data/*\_checks.txt}) & 295 of 295 pass \\
unit tests of the program & 68 pass \\
independent verification, geometry, Dirac field, modes, cosmology & 43
of 43 pass \\
independent verification, self-interaction, Kohn--Sham, CJT, pairing,
gravity & 37 of 37 pass \\
deliberate errors introduced into the independent scripts & 36 of 36
detected \\
checks of the registry of symbols, of the documents and of the data
files against it & 7 of 7 pass \\
numbers of the running text traced to a statement of a program & 500
numbers in 26 documents, 0 not traced \\
\end{longtable}

\hypertarget{conventions-in-one-place}{%
\section{6. Conventions in one place}\label{conventions-in-one-place}}

\begin{itemize}
\tightlist
\item
  Coordinates \(x0,\dots,x7\): \(x1,x2,x3\) ordinary space, \(x4\)
  ordinary time, \(x0\) a hidden space-like coordinate, \(x5,x6,x7\) the
  deflating coordinates.
\item
  Sign of the metric: space-like components positive, time-like
  components negative.
\item
  The rates \(\mathrm{g55}\), \(\mathrm{g66}\), \(\mathrm{g77}\) are
  written upright and never with a subscript; a component of the metric
  carries the superscript \((4+4)\) or \((7+1)\); a coupling of the
  self-interaction, \(g_0,\dots,g_8\), carries one digit as its
  subscript.
\item
  Units: \(\hbar=c=1\).
\item
  In the output of the programs, ``the main project'' means the part of
  PROJECT spacetimes that treats the flat 4+4 and 7+1 spacetimes; its
  program is used as a library, and several checks compare with its data
  files.
\end{itemize}

\hypertarget{what-is-open}{%
\section{7. What is open}\label{what-is-open}}

Three matters need the author's decision: the form of the
self-interaction on the 4+4 spacetime (two forms are computed; one is
adopted), the role of a bare mass beside the Vorton-Scalar, and the
reading of the couplings of the reference example. They are set out,
with the deviations from the instructions and the log of errors, in the
document P0.

\end{document}
